%=============================================================
% Paper 10: Asplund & Nocke (2006), RES
%=============================================================
\chapter{Asplund \& Nocke (2006)}

\begin{paperbox}{Firm Turnover in Imperfectly Competitive Markets}
\textbf{Authors:} Marcus Asplund (LBS/Columbia) \& Volker Nocke (UPenn)\\[3pt]
\textbf{Journal:} \textit{Review of Economic Studies} 73(2): 295--327\\[3pt]
\textbf{Year:} 2006 \quad\textbf{Field:} Applied microeconomics; Industrial
organisation\\[3pt]
\textbf{Classification:} Structural (theoretical + empirical test)
\end{paperbox}

%------------------------------------------------------------
\section{Research Question}

How do market size and fixed costs determine firm turnover (entry and exit
rates) and the age distribution of firms in imperfectly competitive markets?
Does the price-competition effect generated by larger markets systematically
accelerate firm turnover by raising the exit threshold?

%------------------------------------------------------------
\section{Motivation}

Industries differ substantially in firm turnover and entry/exit rates are
highly positively correlated across industries (Dunne, Roberts, Samuelson
1988). Existing theory (primarily Hopenhayn 1992) explains this in perfectly
competitive markets: entry costs and fixed costs affect turnover by changing
the exit threshold. However, in perfectly competitive models, market size has
\emph{no effect} on firm turnover because the entry cost is independent of the
number of firms. In imperfectly competitive industries, a new mechanism
operates: larger markets $\to$ more firms $\to$ lower price-cost margins $\to$
higher exit threshold $\to$ faster turnover. Asplund and Nocke formalise and
test this \textbf{price-competition effect}.

%------------------------------------------------------------
\section{Main Contribution}

\begin{enumerate}
  \item \textbf{Theoretical:} Shows (under Assumption~2: price competition
    effect) that firm turnover is \emph{increasing in market size}---a novel
    prediction unavailable from Hopenhayn's competitive benchmark.
  \item \textbf{Empirical:} Tests the prediction using data on hair salons in
    Sweden across municipalities of different sizes. FOSD tests and regressions
    confirm that firms are younger in larger markets and in markets with higher
    fixed costs.
  \item \textbf{Comparative dynamics:} Derives the full set of signed
    comparative dynamics results for entry costs, fixed costs, and market size.
\end{enumerate}

%------------------------------------------------------------
\section{Relationship to the Literature}

Extends Hopenhayn (1992) to imperfect competition. Differs from Melitz (2003)
in that Dixit-Stiglitz monopolistic competition \emph{does not} satisfy
Assumption~2 (constant markup $\Rightarrow$ no price competition effect;
turnover is independent of market size in Melitz). Standard Cournot and linear
demand models \emph{do} satisfy Assumption~2. Assumption~2 is related to but
finer than Boone's (2000) ``toughness of competition'' analysis.

%------------------------------------------------------------
\section{Theoretical Framework and Economic Mechanism}

\subsection*{Model Structure}

Discrete time, infinite horizon. Firm efficiency $c\in[0,1]$ (lower $c$ =
more efficient). Efficiency evolves as a mixture: with probability $\alpha$,
$c$ stays the same; with probability $1-\alpha$, a new $c$ is drawn from $G(\cdot)$.
Sunk entry cost $\varepsilon>0$; fixed per-period production cost $\phi>0$.

\subsection*{Reduced-Form Profit Function}

Firm gross profit (net of fixed costs) when the measure of active firms is $\mu$:
\begin{equation}
  S\,\pi(c;\mu)\geq 0,
  \label{eq:an_profit}
\end{equation}
where $S$ is market size. Assumptions on $\pi$:
\begin{itemize}
  \item \textbf{(MON):} $\pi(\cdot;\mu)$ strictly decreasing in $c$ on
    $[0,c(\mu))$; zero for $c\geq c(\mu)$.
  \item \textbf{(DOM):} FOSD ordering of firm distributions is preserved by
    the profit ordering.
  \item \textbf{(ORD):} The ordering $(M,\succeq)$ is complete.
  \item \textbf{(CON):} $\pi(c;\mu)$ is continuous.
  \item \textbf{Assumption 1 (market share effect):} For $\mu'\succ\mu$,
    $\pi(c;\mu)-\pi(c;\mu')$ is strictly decreasing in $c$---efficient firms
    suffer larger \emph{absolute} profit reductions from intensified competition.
  \item \textbf{Assumption 2 (price competition effect):} For $\mu'\succ\mu$,
    $\pi(c;\mu')/\pi(c;\mu)$ is strictly decreasing in $c$---less efficient
    firms suffer larger \emph{percentage} profit reductions.
\end{itemize}

Assumption~2 is the key new ingredient. It holds for standard Cournot and
linear demand models but \emph{fails} for Dixit-Stiglitz (constant markup).

\subsection*{Stationary Equilibrium Value Function}

\begin{equation}
  V(c) = \max\bigl\{0,\;\bar{V}(c)\bigr\}, \qquad
  \bar{V}(c) = \bigl[S\pi(c;\mu)-\phi\bigr]
    + \delta\!\left[\alpha V(c)+(1-\alpha)\!\int_0^1 V(z)\,G(dz)\right].
  \label{eq:an_value}
\end{equation}

\subsection*{Entry and Exit Conditions}

In stationary equilibrium with positive entry and exit:
\begin{align}
  &\text{(X)}\quad S\pi(c^*;\mu)-\phi+\delta(1-\alpha)\varepsilon=0
    &&\text{(optimal exit at threshold }c^*\text{)} \\
  &\text{(E)}\quad \int_0^{c^*}\!\bigl[S\pi(c;\mu)-\phi+\delta(1-\alpha)\varepsilon\bigr]G(dc)
    = (1-\alpha\delta)\varepsilon
    &&\text{(free entry)}
\end{align}

The firm at threshold $c^*$ earns negative current net profit
[$S\pi(c^*;\mu)-\phi<0$] but stays because of the option value of getting
a new efficiency draw: $\delta(1-\alpha)\varepsilon>0$.

\subsection*{Central Comparative Dynamics Result (Proposition~4)}

Under Assumption~2, an increase in market size $S$:
\begin{enumerate}
  \item Raises $c^*$ (higher exit threshold---less efficient firms exit).
  \item Increases the rate of firm turnover.
  \item Shifts the age distribution of firms toward younger firms in the
    sense of first-order stochastic dominance.
\end{enumerate}

\textbf{Mechanism:} Larger $S$ $\to$ more entrants $\to$ price-cost margins
fall $\to$ percentage profit reduction is larger for inefficient firms
(Assumption~2) $\to$ their value falls below the exit threshold $\to$ they
exit $\to$ marginal surviving firm must be more efficient $\to$ $c^*$ rises
$\to$ faster turnover $\to$ younger age distribution.

%------------------------------------------------------------
\section{Data}

\begin{itemize}
  \item \textbf{Swedish hair salons (1993--2001):} Age distribution across
    Swedish municipalities of varying population sizes.
  \item \textbf{Advantages:} (i) Single industry $\Rightarrow$ cross-industry
    shock volatility held constant. (ii) Local geographically-distinct markets.
    (iii) Fixed costs vary observably across municipalities.
  \item \textbf{Tests:} Non-parametric FOSD tests of age distributions;
    OLS and quantile regressions of firm age on log population and fixed cost
    proxies.
\end{itemize}

%------------------------------------------------------------
\section{Empirical Strategy}

\begin{enumerate}
  \item For pairs of municipalities of different sizes, test whether the CDF
    of firm ages in the larger market is above (younger) the CDF in the
    smaller market (FOSD test).
  \item Regress firm age on log population (market size) and observable fixed
    cost proxies in OLS and quantile regressions.
  \item Exploit municipality-specific variation in fixed costs to test
    Proposition~5 (higher fixed costs $\Rightarrow$ younger firms).
\end{enumerate}

%------------------------------------------------------------
\section{Identification Strategy}

The key identifying variation is \emph{within-industry, cross-market variation
in market size and fixed costs}. By studying a single industry across many
geographically distinct local markets, the authors eliminate the main confound
in cross-industry tests (cross-industry variation in shock volatility). The
identifying assumption is that local market size (population) is exogenous to
the age distribution of hair salons.

\textbf{Limitation:} No instrumental variable for market size; potential
reverse causality if cities with younger firms attract more consumers.

%------------------------------------------------------------
\section{Main Results}

\subsection*{Theoretical Comparative Dynamics}

\begin{center}
\begin{tabular}{lcc}
\toprule
\textbf{Parameter change} & \textbf{Effect on $c^*$} &
  \textbf{Effect on turnover} \\
\midrule
$\uparrow$ Entry cost $\varepsilon$ & $\downarrow$ & $\downarrow$ \\
$\uparrow$ Market size $S$ & $\uparrow$ & $\uparrow$ \quad (Assumption~2) \\
$\uparrow$ Fixed cost $\phi$ & $\uparrow$ & $\uparrow$ \quad (Assumption~2) \\
\bottomrule
\end{tabular}
\end{center}

\subsection*{Empirical Results (Hair Salons)}

\begin{itemize}
  \item Non-parametric FOSD tests: the age distribution in larger
    municipalities FOSD-dominates that in smaller municipalities---firms are
    \emph{systematically younger} in larger markets. $\checkmark$
  \item OLS: log market size has a negative, statistically significant
    coefficient on average firm age. $\checkmark$
  \item Higher fixed cost proxies: associated with younger firms,
    consistent with Proposition~5. $\checkmark$
  \item Quantile regressions: the effect is pervasive across the full
    age distribution, not driven by outliers.
\end{itemize}

%------------------------------------------------------------
\section{Economic Magnitude and Interpretation}

An increase in market size of one standard deviation reduces average firm age
by a meaningful amount. The finding that higher fixed costs are associated
with \emph{younger} firms (more turnover) is particularly counterintuitive:
it implies that regulations or compliance costs that raise the per-period cost
of operating do not necessarily reduce dynamism---they can \emph{increase}
turnover by forcing out low-productivity incumbents.

%------------------------------------------------------------
\section{Mechanisms}

The price-competition effect operates as:
\begin{enumerate}
  \item Larger market $\to$ free entry drives in more firms.
  \item Price-cost margins fall for all active firms (competition intensifies).
  \item Assumption~2 (price effect): percentage decline in profits is larger
    for less efficient (higher $c$) firms, because price falls are
    proportionally more damaging when margins are already thin.
  \item High-cost incumbents cross the exit threshold first $\to$ they exit.
  \item Marginal surviving firm must be more efficient ($c^*$ rises).
  \item More efficient survivors have shorter expected lifetimes for any
    given entrant $\to$ higher turnover $\to$ younger age distribution.
\end{enumerate}

%------------------------------------------------------------
\section{Robustness Checks}

\begin{itemize}
  \item Main predictions verified under Cournot competition (Appendix);
    Assumptions~1 and~2 hold in this canonical setting.
  \item Results extended to the linear demand model with perceived quality
    differences (isomorphic representation).
  \item FOSD tests are non-parametric---no distributional assumptions.
  \item Quantile regressions confirm results not driven by outliers.
  \item Robustness to alternative fixed cost proxies.
  \item The Dixit-Stiglitz failure of Assumption~2 is explicitly demonstrated
    and discussed as a limitation for trade applications.
\end{itemize}

%------------------------------------------------------------
\section{Limitations}

\begin{enumerate}
  \item \textbf{No IV for market size:} Potential reverse causality.
  \item \textbf{Single industry:} Hair salons are low-tech/labour-intensive;
    generalisability to capital-intensive or high-tech sectors unclear.
  \item \textbf{Dixit-Stiglitz incompatibility:} Most widely used trade
    model (Melitz 2003) does not satisfy Assumption~2.
  \item \textbf{Stylised efficiency process:} Mixture of persistence and
    random redraw; more realistic AR(1) processes may yield different
    quantitative predictions.
  \item \textbf{No worker side:} Welfare effects of higher turnover on
    workers are outside the model.
\end{enumerate}

%------------------------------------------------------------
\section{Policy Implications}

\begin{itemize}
  \item \textbf{Entry regulation:} Lower entry costs increase turnover and
    efficiency (as in Hopenhayn 1992).
  \item \textbf{Fixed production costs:} Higher fixed costs (from compliance
    burdens) increase turnover by forcing exit of less productive firms---
    a counterintuitive policy implication.
  \item \textbf{Economic integration:} Larger integrated markets (EU
    integration, free trade areas) generate more intense price competition,
    raise exit thresholds, and improve aggregate productivity---but at the
    cost of incumbent firm stability.
\end{itemize}

%------------------------------------------------------------
\section{Overall Assessment}

Asplund and Nocke (2006) is an elegant paper that cleanly identifies a new
mechanism (the price-competition effect) not present in Hopenhayn's competitive
benchmark and provides a consistent theory and empirical test. The within-
industry cross-market design is imaginative and appropriate. The main weakness
is the lack of a causal identification strategy for market size and the
restriction to a single low-tech industry. The finding that Dixit-Stiglitz
(most widely used competition model in trade) does not generate the price
competition effect is an important limitation.

\textbf{Quality: Good. Clean theory, appropriate empirical test, but
narrow empirical scope.}

%------------------------------------------------------------
\section{Relevance to My Research}

Relevant for research on firm dynamics in Canada at the RDC. Using LEAP or T2
corporate tax data, one could test whether firm turnover is higher in larger
Canadian metropolitan areas (Toronto vs.\ Prince George), controlling for
industry, consistent with the market size $\to$ turnover prediction.

%------------------------------------------------------------
\section{Potential Research Extensions}

\begin{itemize}
  \item Test the market size $\to$ turnover prediction for Canadian service
    industries (restaurants, retail) using LEAP/LEED data across municipalities.
  \item Examine whether CETA/CPTPP increased firm turnover in Canadian
    import-competing industries through the price-competition channel.
\end{itemize}

%------------------------------------------------------------
\section{Research Gaps}

\begin{itemize}
  \item The aggregate productivity consequences of market-size-driven turnover
    are not quantified---a Foster-Haltiwanger-Krizan (2001) decomposition
    of productivity growth across markets of different sizes would be the
    natural next step.
  \item No worker side: welfare effects of higher firm turnover on workers
    (displacement, wage changes) are outside the model.
\end{itemize}

%------------------------------------------------------------
\section{Methodological Lessons}

\begin{enumerate}
  \item Testing theory within a \emph{single industry} across geographic
    markets is a powerful design eliminating the main confound (cross-industry
    shock volatility) in cross-industry tests.
  \item Non-parametric FOSD tests of age distributions are more convincing
    than tests based on a single summary statistic (mean age), as they test
    the full distributional prediction.
  \item The Assumption~1/2 taxonomy (market share effect vs.\ price
    competition effect) is a useful organising framework for comparing models
    of imperfect competition on their firm dynamics implications.
\end{enumerate}

%------------------------------------------------------------
\section*{Five-Sentence Summary}

Asplund and Nocke (2006) develop a stochastic dynamic model of a monopolistically
competitive industry showing that firm turnover is increasing in market size
through a price-competition effect: larger markets support more firms, which
intensifies competition, and the percentage profit reduction is larger for
less efficient (higher-cost) firms (Assumption~2), raising the exit threshold
and accelerating turnover. This prediction---absent from Hopenhayn's (1992)
competitive benchmark and inconsistent with the Dixit-Stiglitz model---is
tested using data on Swedish hair salons across municipalities of varying
sizes, where within-industry geographic variation eliminates the main confound
in cross-industry tests. Non-parametric FOSD tests and OLS regressions confirm
that larger markets have age distributions shifted toward younger firms, and
markets with higher fixed costs (Proposition~5) similarly have younger firms.
The model delivers clean comparative dynamics: higher entry costs reduce
turnover by protecting incumbents, while higher fixed production costs increase
turnover through a selection effect. The paper establishes that market size
is a first-order determinant of firm dynamics under imperfect competition---
a mechanism missing from competitive industry models and from the widely used
Dixit-Stiglitz/Melitz trade framework.

\vspace{6pt}
\begin{quickref}
\begin{tabularx}{\linewidth}{@{}lX@{}}
\textbf{Key contribution} &
  Novel theoretical prediction that firm turnover is increasing in market
  size under imperfect competition (price-competition effect); first paper
  to embed this in a stochastic dynamic equilibrium and test it empirically
  with a within-industry cross-market design. \\[4pt]
\textbf{Key weakness} &
  No IV for market size (potential reverse causality); limited to a single
  low-tech industry (hair salons in Sweden); central result does not hold
  under Dixit-Stiglitz (most common trade model). \\[4pt]
\textbf{Most interesting gap} &
  The aggregate productivity implications of market-size-driven turnover
  are not quantified; a Foster-Haltiwanger-Krizan (2001) decomposition of
  productivity growth across markets of different sizes would complete the
  welfare analysis. \\[4pt]
\textbf{Publishable extension} &
  Test the Asplund-Nocke market size $\to$ firm turnover prediction for
  Canadian service industries using LEAP data across Census Metropolitan
  Areas; study whether CETA/CPTPP increased turnover in import-competing
  industries through the price-competition channel.
\end{tabularx}
\end{quickref}
